% ==========================================================================
% GENERATED by scripts/split_paper.py. Do not edit this file.
%
% Source: imported_papers/Llama2_benchmark_sds/llm_at_home.tex
% Part 3: Benchmarking Performance and Quality
%
% Edits belong upstream, in the published paper. To change what the
% chapter includes, change imports/llmbench.toml and re-run
% the splitter.
% ==========================================================================
\section{Benchmarking Performance and Quality}
\label{sec:bench}
We present benchmarks for two distinct use cases: First, we assess model inference by employing a pre-trained model. Second, we examine fine-tuning on custom data. Model quantization and LoRA are the most promising, impactful, and universally applicable among the plethora of techniques introduced above. Therefore, we chose them for benchmarking. In both scenarios, we conduct a comparative analysis between a given model and the same model incorporating additional quantization techniques. In the inference scenario, we benchmark quantization techniques, while in the fine-tuning scenario, we test multiple low-rank approximations. One focus lies on evaluating VRAM usage, speed, and model performance, as these metrics are pivotal for the effective deployment of models in resource-constrained environments. The second focus lies on gauging the quality of the generated conversations. Since it is the most capable openly available LLM, we use LlamaV2 \cite{touvron_llama_2023} as the basis for all of our experiments. We use its 7-billion parameter variant -- larger versions require a substantial GPU infrastructure, even when paired with the techniques introduced above. 

\subsection{Benchmarking methodology}
\label{sec:bench_intro}
% \subsubsection{Methods to be benchmarked}

% We benchmark two scenarios inference where an appropriate open-source LLM exists and the challenge is only to deploy it locally and fine-tuning where an existing LMM has to be further trained on a new dataset.
% \\ \emph{Inference:}
% For inference-only experiments, we use LlamaV2-Chat which has been extensively tuned to be a chat based assistant before release. As a baseline, we run inference on LlamaV2-Chat without any quantization (no quant), we then run inference using a fixed quantization as implemented in the package \emph{bitnbytes} \cite{bnb}. Rounding out the lineup are two adaptive quantization methods \emph{awq} \cite{awq} and \emph{gptq} \cite{gptq}.
% \\ \emph{Fine-tuning:}
% 


\subsubsection{Resource usage and speed metrics}
For inference, we systematically report metrics acquired during the quantitative evaluation (c.f. Section \ref{sec:quantitative_evaluation}) on the MT benchmark. This includes VRAM usage (the most constraining factor for the seamless execution of state-of-the-art models with billions of parameters), generated tokens per second, and overall power consumption in watt-hours (Wh).
For fine-tuning, we provide metrics obtained during the re-training, including VRAM usage, processed batches per second, and overall power consumption. 
For simplicity, we constrain the models to a batch size of $1$, even though the aforementioned methods might reduce VRAM requirements to allow an increase in batch size, contingent on available computational resources.

The experiments are conducted using a single Nvidia A-100 $40$GB GPU. The selection of this GPU was based on its ample VRAM, allowing the execution of all model versions (except direct fine-tuning which is out of the scope of this work), both with and without quantization, facilitating a fair comparison. We also report results generated using a Nvidia T4 $16$GB GPU to demonstrate the feasibility of running and fine-tuning models on constrained hardware (``at home''). It is important to note that, due to the memory limitations of the Nvidia T4 $16$GB GPU, only models employing quantization methods could be executed on it.

\subsubsection{Quantitative Evaluation}\label{sec:quantitative_evaluation}
We employ the MT-benchmark that contains $80$ high-quality multi-turn questions described in \cite{zheng_judging_2023}. They are manually designed to evaluate the models' proficiency in multi-turn conversation and instruction-following capabilities. The questions are organized into eight categories, each holding ten challenging questions: writing, roleplay, extraction, reasoning, math, coding, STEM knowledge, and humanities/social science knowledge. Each multi-turn question comprises an initial plus a subsequent inquiry.

The multi-turn questions are fed into all model versions, and the resulting responses are stored. Subsequently, OpenAI's GPT-4 model is utilized to assess response quality by assigning a score ranging from $1$ to $10$ to each answer.
To obtain the score, we follow the prompting scheme of Zheng et al. \cite{zheng_judging_2023}, asking GPT-4 to evaluate the model's response for various factors such as ``relevance'', ``accuracy'', and ``depth'' before providing an overall score.
The average score per category is then computed to provide a comprehensive evaluation of the model's performance across different categories. In addition to this single-answer evaluation, a pairwise comparison is conducted between the baseline model and a second model version that has been tuned for efficiency. In this process, the responses from both models are presented to the auxiliary LLM, GPT-4, which is then tasked with determining the superior model or declaring a tie.
Both procedures adhere to the methodology proposed by Zheng et al. \cite{zheng_judging_2023}.

\subsubsection{Qualitative Evaluation}
\label{sec:quali_eval}
As automatic evaluation of generative language models is still in its infancy, it is important to also qualitatively inspect outputs to assess the quality of a model. We systematically do this by prompting our models with the same catalog of five questions sourced from MT-bench. Then, eight individual human judges rate all the outputs on a scale from $1$ (useless) to $5$ (perfect answer). Finally, we average across the marks given by each judge to obtain an objective score. We use the following five multi-turn questions to assess the ability of the models on a broad range of topics:

\small
\begin{enumerate}
     \item[\textbf{Q1}] (Writing) Round 1: \emph{Compose an engaging travel blog post about a recent trip to Hawaii, highlighting cultural experiences and must-see attractions.} \\
     Round 2: \emph{Rewrite your previous response. Start every sentence with the letter A.}
     
     \item[\textbf{Q2}] (Reasoning) Round 1: \emph{Thomas is very healthy, but he has to go to the hospital every day. What could be the reasons?} \\
     Round 2: \emph{Can you explain why the above question is interesting?}

     \item[\textbf{Q3}] (Math) Round 1: \emph{Benjamin went to a bookstore and purchased a variety of books. He bought $5$ copies of a sci-fi novel, each priced at \$$20$, $3$ copies of a history book priced at \$$30$ each, and $2$ copies of a philosophy book for \$$45$ each. What was the total cost of his purchases?} \\
     Round 2: \emph{Suppose Benjamin decides to sell each of these books at a $25$\% markup from the price he purchased them. What would be his total revenue if he sold all the books he bought?}

     \item[\textbf{Q4}] (Stem) Round 1: \emph{Photosynthesis is a vital process for life on Earth. Could you outline the two main stages of photosynthesis, including where they take place within the chloroplast, and the primary inputs and outputs for each stage?} \\
     Round 2: \emph{How much energy can a tree produce through photosynthesis in its lifetime? Please provide an estimate using actual numerical values and thoroughly explain your thought process step-by-step.}

     \item[\textbf{Q5}] (Humanities) Round 1: \emph{Which methods did Socrates employ to challenge the prevailing thoughts of his time?} \\
     Round 2: \emph{Let's bring Socrates to the modern world. Generate a conversation between Socrates and Bill Gates to debate on generative AI for education.}
\end{enumerate}
\normalsize

\subsection{Results for Efficient Inference}
\label{sec:bench_inference}

In the subsequent analysis, we assess different quantization techniques with a focus on computational, quantitative, and qualitative performance. Out of the numerous open-source LLMs \cite{zhao_survey_2023}, we selected``LlamaV2-7B'' and its dialogue-fine-tuned version ``LlamaV2-7B-Chat'' as the basis for this evaluation. This choice is based on the models' good performance with a manageable, popular parameter count of ``just'' $7$B \cite{touvron_llama_2023}.

\begin{table*}
\caption{VRAM usage, tokens processed per second, and the overall power consumption during inference on MT-bench. The model used is LlamaV2-Chat, optionally combined with bits-and-bytes (BnB), nested BnB, AWQ, and GPTQ quantization techniques.}
\label{tab:inference_metrics}
\centering
\begin{tabular}{lccccc}
                & VRAM usage [GB] & \multicolumn{2}{c}{avg. tokens/s} & \multicolumn{2}{c}{power consumption [Wh]}  \\  \hline
GPU                        & A100      & A100  & T4  & A100  & T4 \\ \hline
LlamaV2-7B-Chat            & $33.80$ & $36.06$ & out of memory & $92.39$ & out of memory \\
LlamaV2-7B-Chat BnB        & $8.40$ & $25.45$ & $16.07$ & $63.53$ & $57.09$ \\
%LlamaV2-7B-Chat nested BnB& 8.00 & 19.78 & 16.06 & 69.15 & 60.90 \\
LlamaV2-7B-Chat AWQ        & $9.10$ & $86.66$ & $26.32$ & $33.72$ & $38.96$ \\
LlamaV2-7B-Chat GPTQ       & $9.02$ & $34.20$ & $28.33$ & $52.56$ & $32.28$ \\ \hline
\end{tabular}
\end{table*}

\emph{Resource usage and speed metrics.}
Table \ref{tab:inference_metrics} presents an overview of the computational requirements of the LlamaV2-7B-Chat model and its variants employing quantization techniques during the processing of the MT-bench dataset. The first column shows the VRAM usage. It is important to acknowledge that VRAM utilization may vary based on the sampled model response that is reused as conversational context in subsequent follow-up questions.
The adoption of quantization techniques enables the execution of the model on GPUs with constrained VRAM, such as the Nvidia T4, as the VRAM requirements reduce from $33.80$GB to below $10$GB. Alternatively, in scenarios where sufficient memory is available, this facilitates the utilization of batch sizes that are three to four times larger, thereby enabling the concurrent processing of a higher volume of requests.

When examining the tokens processed per second using the A-100 GPU, it becomes evident that quantization via BnB ($25.45$ tokens/s) exhibits significantly slower performance compared to the baseline without quantization ($36.06$ tokens/s).
Thus, squeezing parameters and activations into lower precision datatypes reduces memory requirements on the cost of speed.
In contrast, more advanced algorithms such as GPTQ and AWQ can even preserve or improve the throughput while having a reduced memory footprint:
GPTQ quantization achieves a throughput of $34.2$ tokens/s, nearly matching the baseline. Remarkably, AWQ quantization attains an impressive throughput of $86.66$ tokens/s, approximately $2.4$ times faster than the baseline.
Unfortunately, this enhanced performance does not scale to the T4 GPU due to this GPU's reliance on the older Turing architecture. The T4 GPU lacks the novel CUDA kernels of the Ampere architecture, which is extensively utilized by AWQ. Consequently, AWQ on the T4 GPU achieves performance similar to GPTQ.
However, as previously elucidated, the quantization methods enable larger batch sizes and the concurrent processing of more tokens, thereby enhancing overall throughput. Thus, even slower methods like BnB could outperform the baseline when quantization is employed to increase the batch size.

Quantization techniques also allow for substantially diminishing power consumption. This is attributed to the utilization of lower-precision datatypes for representing model parameters and activations, resulting in enhanced computational efficiency. On the A-100 GPU, the standard model consumes $92.39$Wh, whereas BnB has a power consumption of $63.53$Wh, and nested BnB consumes $69.15$Wh. GPTQ and AWQ employ similar datatypes as BnB, but due to their higher throughput, these methods achieve an even lower total power consumption of $33.72$Wh (AWQ) and $52.56$Wh (GPTQ), respectively.
Unfortunately, due to the aforementioned reasons, AWQ struggles to transfer its high efficiency to the older GPU.

To summarise, from a purely computational perspective, AWQ quantization is preferable if the GPU is equipped with compatible CUDA drivers; otherwise, GPTQ is recommended. These approaches demand significantly less VRAM compared to the baseline model while maintaining or increasing processing speed, even when the batch size remains unchanged. Unfortunately, AWQ and GPTQ cannot be used for model training, making BnB the preferred method for quantization while learning model weights.

\begin{figure}[!t]
\centering
\includegraphics[width=0.5\textwidth]{imgs/single_2}
\caption{Single score evaluation results for LlamaV2-7B-Chat  with different quantization techniques and the state-of-the-art GPT-4 for comparison.}
\label{fig:single_2}
\end{figure}
%
\emph{Quantitative evaluation.}
%
To quantitatively assess the influence of quantization on LlamaV2-7B-Chat's output quality, we utilize the single-score judgment based on MT-bench. Additionally, we include evaluations of LlamaV2-7B as a lower bound and GPT-4 as an upper bound. Please note that GPT-4, when used as a judge, favors its own answers with a $10\%$ bias in pair-comparisons, suggesting that it rates itself also slightly better than other models in single answer grading \cite{zheng_judging_2023}. LlamaV2-7B is used as the lower bound due to its identical architecture but lack of fine-tuning for dialogue, a requirement for MT-bench. 
The findings are presented in Figure \ref{fig:single_2}.
As expected, the LlamaV2-7B-Chat model and its quantized versions are noticeably better than the lower bound and worse than the upper bound. 
There is only a marginal performance difference between the baseline and the quantized models, suggesting that the application of quantization techniques neither degrades nor enhances performance.


\emph{Qualitative evaluation.}
Table \ref{tab:qual_inference} shows the average grade for the answers of all the quantized versions of LlamaV2-Chat to our five selected questions (see Section \ref{sec:quali_eval}). 
Notably, the non-quantized model achieves the highest average grade of $3.35$, surpassing GPTQ with an average grade of $3.1$ and BnB with an average grade of $2.9$. In contrast, AWQ received a lower average rating of $1.95$ compared to the other quantization methods. This discrepancy suggests that humans do not favor AWQ as a quantization technique, which is in contrast to the outcomes of the quantitative analysis, where no performance gap could be observed. Our results indicate that quantization techniques improve computational efficiency but degrade the overall dialogue quality.

\begin{table}
\centering
\caption{Average (across $8$ human raters) usefulness evaluation ($1$=useless, $5$=perfect) of answers from quantized versions of LlamaV2-Chat}
\label{tab:qual_inference}
\begin{tabular}{lllllll}
         & Q1   & Q2   & Q3   & Q4   & Q5   & Average \\ \hline
no quantization & $3.75$ & $1.25$ & $4.5$  & $4$    & $3.25$ & $3.35$    \\
BnB      & $2.25$ & $2$    & $2.75$ & $3.5$  & $4$    & $2.9$     \\
AWQ      & $1.5$  & $1.5$  & $1.5$  & $2.5$  & $2.75$ & $1.95$    \\
GPTQ     & $1.25$ & $1.75$ & $4.75$ & $3.75$ & $4$    & $3.1$     \\ \hline
\end{tabular}
\end{table}

\subsection{Results for Efficient Model Fine-Tuning}
\label{sec:bench_ft}
When there is no task-appropriate LLM available, it becomes necessary to fine-tune an LLM on a new dataset. To evaluate the effect of efficiency techniques for LLM fine-tuning, we fine-tune a LlamaV2 base model on the dataset that was designed for the training of Guancano \cite{qlora} and that is a subset of the OpenAssistant \cite{kopf_openassistant_2023} dataset. Thus, we train the model to behave like a chat-based assistant rather than a general generative model. While this is of no 
practical importance as a new model due to the existence of LlamaV2-Chat, it allows us to keep a consistent evaluation scheme across all experiments for this benchmarking.
Fine-tuning a LlamaV2 model without a low-rank adapter is only possible with very large computational resources (at least $80$GB VRAM on a single chip), therefore we consider the widely used LoRA \cite{lora} method as our baseline. Additionally, we fine-tune models using AdaLoRA \cite{adalora} and LoKr \cite{LoKr}. We also test quantized versions of LoRA (Q-LoRA) and AdaLoRA (Q-AdaLoRA), which provide a further reduction in VRAM requirement.
We train our models with a constant learning rate of $0.0001$ and a weight decay of $0$. We use a batch size of $1$ and $8$ gradient accumulation steps. We cap the gradient norm at $0.3$ and train for $15,000$ steps. LoKr has a relatively heavy VRAM load, therefore we have to reduce the maximum sequence length to $400$ tokens instead of the default of $500$, even on a $40$GB GPU. To enable fair comparisons, we keep this training configuration (including max sequence length) consistent across all methods. The full training details can be found in the public WandB logs linked here: \url{https://github.com/tuggeluk/LMM_at_home}. We have opted for this relatively short fine-tuning due to our limited computational resources. This setup allows us to compare the different low-rank adaptation methods, but the resulting models are not expected to be competitive with the highly tuned LlamaV2-Chat.  

\emph{Resource usage and speed metrics.} 
\begin{table*}
\caption{VRAM usage, tokens processed per second, and overall power consumption during fine-tuning. The base model used is LlamaV2-7B, fine-tuned with low-rank approximators LoRA, AdaLoRA, LoKr, quantized LoRA, and quantized AdaLoRA.}
\label{tab:finetune_metrics}
\centering
\begin{tabular}{lcccccc}
                & VRAM usage [GB]  & \multicolumn{2}{c}{steps/s} & \multicolumn{2}{c}{power usage [Wh]}  \\  \hline
GPU             &      & A100  & T4  & A100  & T4 \\ \hline
LoRA            &  $33.4$  & $0.217$ & out of memory & $6,641$ & out of memory \\   
AdaLoRA         &  $32.7$ & $0.203$ & out of memory &  $6,560$ & out of memory \\   
LoKr            &  $39.5$ & $0.15$  & out of memory &  $10,080$ & out of memory \\
Q-LoRA          &  $13.5$ & $0.178$ & $0.04$ &  $7,990$ & $7,402$ \\ 
Q-AdaLoRA       &  $12.7$ & $0.168$ & $0.04$ &  $7,936$ & $7,843$ \\ \hline            
\end{tabular}
\end{table*}
%
Table \ref{tab:finetune_metrics} shows an overview of the resources consumed by the different fine-tuning methods for the A100 and T4 GPUs. Among the non-quantized adapters, LoKr requires a higher amount of memory than LoRA and AdaLoRA. Quantizing the adapter slashes the memory required for fine-tuning by almost two-thirds, allowing the use of T4 chips that only have $16$GB of memory. The processing speed does not deviate too sharply between the methods with the fastest processing (LoRA: $0.217$ steps/s) and the slowest processing (LoKr: $0.15$ steps/s) on an A-100. Quantizing causes a slowdown of roughly $20\%$ in both cases, comparable to the BnB quantization during inference. The T4 is much slower at $0.04$ steps/s for both methods,  yielding an increase in training time of more than four-fold. The overall power consumption on the A-100 is highly correlated with the training time, while the reduced memory usage of the quantized methods has a slim effect on the energy consumption, making the faster, unquantized training a more energy-efficient option. While the T4 has a much lower peak power consumption, this is largely compensated by longer training time, leading to similar energy usage for both GPU models.


\begin{figure}[!t]
\centering
\includegraphics[width=0.5\textwidth]{imgs/pair_2}
\caption{Pair-wise comparison results of LlamaV2-7B vs. fine-tuned models.}
\label{fig:pair}
\end{figure}
%
\emph{Quantitative evaluation.}
To assess the impact of efficiency-optimized fine-tuning on chat quality, we conduct a pair-wise comparison using MT-bench. In this comparison, a judging model assigns scores of $1.0$ if the fine-tuned version is superior to the baseline model, $0.5$ if they perform equally, and $0.0$ if the fine-tuned version is inferior. The results, depicted in Figure \ref{fig:pair}, reveal that our fine-tuned models consistently outperform the baseline, consisting of a stock LlamaV2, across all categories. This underscores the effectiveness of our fine-tuning.
Notably, the fine-tuned models exhibit especially superior performance in the categories of writing, roleplay, extraction, STEM, and humanities, with average scores hovering around $1.0$, indicating their consistent superiority over the baseline in nearly every instance.

\begin{figure}[!t]
\centering
\includegraphics[width=0.5\textwidth]{imgs/single_1}
\caption{Single score evaluation results for pre-trained LlamaV2-7B (baseline) and our fine-tuned versions.}
\label{fig:single_1}
\end{figure}
%
In addition to conducting pairwise comparisons, we present outcomes based on single-score assessments using MT-bench. The outcomes are illustrated in Figure \ref{fig:single_1}. The results from our fine-tuned models follow a similar pattern as the results from the quantized LlamaV2-Chat version used for inference (cf. Figure \ref{fig:single_2}), i.e., work better in specific categories than others. However, overall, they perform worse than the models used during inference.
This discrepancy arises from our use of a simplified and short fine-tuning approach, wherein we refrain from manually evaluating multiple checkpoints to identify the optimal model version. Nonetheless, this validation affirms that our results are comparable to the open-source fine-tuned versions.

When comparing the results from the fine-tuned versions, only subtle distinctions are evident. For instance, LoKr exhibits a slightly inferior performance in tasks related to humanities and STEM compared to the other models. Conversely, LoKr, along with AdaLora, surpasses other models in tasks involving information extraction.
Moreover, quantization methods appear to have negligible effects on the overall quality of the model. The mean rating for LoRA stands at $4.58$, and its quantized counterpart, Q-LoRA, is rated at $4.63$ ($+0.05$). Similarly, AdaLoRA achieves an average score of $4.96$, and its quantized variant, Q-AdaLoRA, is rated at $4.60$ ($-0.36$). These variances are minimal, especially considering the scale of ratings ranging from one to ten.


%\begin{table*}[]
%\caption{The pair-wise comparison results of the fine-tuned LlamaV2-7B Lora vs. other quantization methods. The scores are the average over $10$ multi-turn question results, where $0.0$ means that the first model is always better, $0.5$ is a tie, and $1.0$ means the second model is always better.}
%\label{tab:fine-tab:fine-tuning_pair_scores}
%\centering
%\begin{tabular}{lcccccccc}
%                & Writing & Roleplay & Reasoning & Math & Coding & Extraction & STEM & Humanities \\ \hline
%AdaLora vs. Lora    & 0.45 & 0.58 & 0.48 & 0.50 & 0.50 & 0.60 & 0.30 & 0.35\\
%LoKr vs. Lora       & 0.40 & 0.40 & 0.65 & 0.45 & 0.45 & 0.60 & 0.18 & 0.18\\
%Q-Lora vs. Lora     & 0.43 & 0.55 & 0.50 & 0.50 & 0.48 & 0.38 & 0.53 & 0.48 \\
%Q-AdaLora vs. Lora  & 0.52 & 0.48 & 0.53 & 0.45 & 0.58 & 0.55 & 0.33 & 0.38\\ \hline
%\end{tabular}
%\end{table*}
%
%
%\begin{figure}[!t]
%\centering
%\includegraphics[width=0.5\textwidth]{imgs/pair_1}
%\caption{The pair-wise comparison results of the fine-tuned LlamaV2-7B Lora vs. other quantization methods.}
%\label{fig:pair}
%\end{figure}


\emph{Qualitative evaluation.}
Table \ref{tbl:qualy_finetuned} shows the average grade for answers of all fine-tuned models to our five selected questions (see Section \ref{sec:quali_eval}). We observe that, on average, LoRA performs the best, with a strong performance gap to AdaLoRA that in turn outperforms LoKr. Interestingly, the human evaluation of the answers does indicate a quite steep drop-off in quality between the quantized and unquantized versions of LoRA and AdaLoRA. This is a further indication that model quantization does not incur just a negligible loss in performance as most automated evaluations suggest.

\begin{table}
\centering
\caption{Average (across $8$ human raters) usefulness evaluation ($1$=useless, $5$=perfect) of answers from fine-tuned LlamaV2s.}
\label{tbl:qualy_finetuned}
\begin{tabular}{lllllll}
         & Q1  & Q2  & Q3  & Q4   & Q5  & Average \\ \hline
LoRA     & $2.5$  & $2.25$ & $0.75$ & $3.75$ & $3.25$ & $2.5$     \\
Q-LoRA    & $2$    & $0.75$ & $2.25$ & $3.25$ & $1.75$ & $2$       \\
AdaLoRA  & $1.5$  & $1.25$ & $2.5$  & $3$    & $2.25$ & $2.1$     \\
Q-AdaLoRA & $1.25$ & $1.75$ & $1.25$ & $1.5$  & $2.25$ & $1.6$     \\
LoKr     & $2.25$ & $1.5$  & $1$    & $0.5$  & $3$    & $1.65$    \\ 
\hline
\end{tabular}
\end{table}


% \subsection{Text classification}
% Quantitatively evaluate trained models on a text classification dataset. (baseline: just train classification layer --> how much does fine-tuning of the LLM help?)
% --> maybe skip this chapter


% \subsection{Conversational model}
% Qualitatively evaluate trained models on a conversational dataset on the same 2-3 conversations

\subsection{Colab Notebooks to ``Run Your LLM at Home''}
\label{sec:notebooks}
%Link to all of our experiments to be replicated on collab --> major selling point that everyone can replicate this at home.
We provide easy-to-use Colab notebooks for all presented methods that work on a T4 GPU (see \url{https://github.com/tuggeluk/LMM_at_home}). Using these notebooks, all quantized inference methods as well as Q-LoRA and Q-AdaLoRA fine-tuning can be easily deployed on a Colab free-tier instance by anyone without any coding.
